% ============================================================
\chapter{Plongements de Mots}
\label{ch:plongements}
% ============================================================

En 2013, Tomas Mikolov et son \'equipe chez Google publient un article qui va transformer le traitement du langage naturel~: \emph{Efficient Estimation of Word Representations in Vector Space}. L'id\'ee est d'entra\^iner un r\'eseau de neurones superficiel \`a pr\'edire les mots \`a partir de leur contexte (ou inversement), et d'utiliser les poids appris comme repr\'esentations vectorielles des mots. Le r\'esultat, Word2Vec, r\'ev\`ele une structure g\'eom\'etrique stupéfiante~: les analogies s\'emantiques deviennent des \emph{op\'erations vectorielles}. Le vecteur \og roi \fg $-$ \og homme \fg $+$ \og femme \fg pointe vers \og reine \fg. La g\'eom\'etrie capture le sens.

Mais Word2Vec n'est pas n\'e de rien. Il s'inscrit dans une longue tradition fond\'ee sur l'\emph{hypoth\`ese distributionnelle} de J.R.~Firth (1957)~: \og on conna\^it un mot par la compagnie qu'il fr\'equente. \fg Ce chapitre retrace le chemin qui m\`ene des matrices de co-occurrence aux plongements denses, en passant par l'analyse s\'emantique latente et GloVe.

\begin{intuition}[title=\textbf{L'hypoth\`ese distributionnelle}]
\og You shall know a word by the company it keeps \fg (J.R. Firth, 1957).
Les plongements de mots (\emph{word embeddings}) sont des repr\'esentations
vectorielles denses qui capturent la similarit\'e s\'emantique : les mots
apparaissant dans des contextes similaires ont des vecteurs proches.
\end{intuition}

% ------------------------------------------------------------
\section{De la repr\'esentation creuse \`a la repr\'esentation dense}
\label{sec:plg:motivation}
% ------------------------------------------------------------

Rappelons qu'une repr\'esentation one-hot encode chaque mot $w_i$ comme un
vecteur $\mathbf{e}_i \in \R^V$ avec $(\mathbf{e}_i)_j = \delta_{ij}$.
Cette repr\'esentation souffre de deux d\'efauts majeurs :
\begin{enumerate}
    \item \textbf{Dimensionnalit\'e} : $V$ peut atteindre $10^5$ \`a $10^6$.
    \item \textbf{Orthogonalit\'e} : $\inner{\mathbf{e}_i}{\mathbf{e}_j} = 0$
          pour $i \neq j$, m\^eme pour des synonymes.
\end{enumerate}

L'objectif des plongements est d'apprendre une fonction
$\phi : \mathcal{V} \to \R^d$ avec $d \ll V$ (typiquement $d \in [100, 300]$)
telle que la proximit\'e g\'eom\'etrique refl\`ete la similarit\'e s\'emantique.

% ------------------------------------------------------------
\section{Word2Vec}
\label{sec:plg:word2vec}
% ------------------------------------------------------------

\subsection{Architecture CBOW}

\begin{definition}[Continuous Bag of Words (CBOW)]
Le mod\`ele CBOW pr\'edit le mot central $w_t$ \`a partir de son contexte
$C_t = \{w_{t-c}, \ldots, w_{t-1}, w_{t+1}, \ldots, w_{t+c}\}$ o\`u $c$ est
la taille de la fen\^etre. L'objectif est de maximiser :
\[
\mathcal{L}_{\text{CBOW}} = \sum_{t=1}^{T} \log P(w_t \mid C_t)
\]
avec :
\[
P(w_t \mid C_t) = \frac{\exp(\inner{\mathbf{v}_{w_t}}{\bar{\mathbf{u}}_{C_t}})}{\sum_{w \in \mathcal{V}} \exp(\inner{\mathbf{v}_w}{\bar{\mathbf{u}}_{C_t}})}
\]
o\`u $\bar{\mathbf{u}}_{C_t} = \frac{1}{2c} \sum_{w \in C_t} \mathbf{u}_w$
est la moyenne des vecteurs de contexte.
\end{definition}

\subsection{Architecture Skip-gram}

\begin{definition}[Skip-gram]
Le mod\`ele Skip-gram pr\'edit les mots du contexte \`a partir du mot central.
L'objectif est de maximiser :
\[
\mathcal{L}_{\text{SG}} = \sum_{t=1}^{T} \sum_{\substack{j=-c \\ j \neq 0}}^{c} \log P(w_{t+j} \mid w_t)
\]
avec :
\[
P(w_o \mid w_i) = \frac{\exp(\inner{\mathbf{v}_{w_o}}{\mathbf{u}_{w_i}})}{\sum_{w \in \mathcal{V}} \exp(\inner{\mathbf{v}_w}{\mathbf{u}_{w_i}})}
\]
\end{definition}

\begin{attention}[title=\textbf{Co\^ut du softmax}]
Le d\'enominateur du softmax n\'ecessite une somme sur tout le vocabulaire $V$,
ce qui est prohibitif. Deux approximations sont couramment utilis\'ees :
l'\'echantillonnage n\'egatif et le softmax hi\'erarchique.
\end{attention}

\subsection{\'Echantillonnage n\'egatif (NEG)}

\begin{definition}[Negative Sampling]
L'objectif avec \'echantillonnage n\'egatif remplace le softmax par :
\[
\mathcal{L}_{\text{NEG}} = \log \sigma(\inner{\mathbf{v}_{w_o}}{\mathbf{u}_{w_i}}) + \sum_{k=1}^{K} \E_{w_k \sim P_n} \left[\log \sigma(-\inner{\mathbf{v}_{w_k}}{\mathbf{u}_{w_i}})\right]
\]
o\`u $K$ est le nombre d'exemples n\'egatifs et $P_n(w) \propto \text{count}(w)^{3/4}$.
\end{definition}

\begin{theoreme}[\'Equivalence NEG et PMI]
\label{thm:neg-pmi}
Levy et Goldberg (2014) ont montr\'e que Skip-gram avec \'echantillonnage
n\'egatif factorise implicitement une matrice de PMI d\'ecal\'ee :
\[
\mathbf{u}_w^\top \mathbf{v}_c = \text{PMI}(w, c) - \log K
\]
o\`u $\text{PMI}(w, c) = \log \frac{P(w, c)}{P(w) P(c)}$ est l'information
mutuelle ponctuelle.
\end{theoreme}

% ------------------------------------------------------------
\section{GloVe}
\label{sec:plg:glove}
% ------------------------------------------------------------

\begin{definition}[GloVe]
Le mod\`ele GloVe (Pennington et al., 2014) apprend les plongements en
minimisant l'erreur quadratique pond\'er\'ee sur la matrice de co-occurrence :
\[
\mathcal{L}_{\text{GloVe}} = \sum_{i,j=1}^{V} f(X_{ij}) \left(\inner{\mathbf{u}_i}{\mathbf{v}_j} + b_i + c_j - \log X_{ij}\right)^2
\]
o\`u $X_{ij}$ est le nombre de co-occurrences, $b_i, c_j$ sont des biais, et
$f$ est une fonction de pond\'eration :
\[
f(x) = \begin{cases}
(x / x_{\max})^\alpha & \text{si } x < x_{\max} \\
1 & \text{sinon}
\end{cases}
\]
avec typiquement $x_{\max} = 100$ et $\alpha = 0.75$.
\end{definition}

\begin{remarque}
GloVe combine les avantages des m\'ethodes globales (factorisation matricielle)
et locales (fen\^etre de contexte). L'objectif GloVe est \'equivalent \`a une
factorisation de la matrice de log-co-occurrence.
\end{remarque}

% ------------------------------------------------------------
\section{FastText}
\label{sec:plg:fasttext}
% ------------------------------------------------------------

\begin{definition}[FastText]
FastText (Bojanowski et al., 2017) enrichit Skip-gram en repr\'esentant chaque
mot comme la somme de ses $n$-grammes de caract\`eres. Pour un mot $w$ avec
l'ensemble de $n$-grammes $\mathcal{G}_w$ :
\[
\mathbf{u}_w = \sum_{g \in \mathcal{G}_w} \mathbf{z}_g
\]
o\`u $\mathbf{z}_g \in \R^d$ est le vecteur du $n$-gramme $g$.
\end{definition}

\begin{bonnepratique}[title=\textbf{Avantages de FastText}]
\begin{itemize}
    \item G\`ere les mots hors vocabulaire (OOV) par d\'ecomposition en $n$-grammes.
    \item Capture la morphologie (pr\'efixes, suffixes).
    \item Particuli\`erement efficace pour les langues morphologiquement riches
          (turc, finnois, arabe).
\end{itemize}
\end{bonnepratique}

% ------------------------------------------------------------
\section{Propri\'et\'es g\'eom\'etriques}
\label{sec:plg:geometrie}
% ------------------------------------------------------------

\subsection{Analogies vectorielles}

L'une des propri\'et\'es les plus remarquables des plongements est leur capacit\'e
\`a capturer des relations s\'emantiques par des op\'erations arithm\'etiques :

\begin{exemple}
\[
\mathbf{v}_{\text{roi}} - \mathbf{v}_{\text{homme}} + \mathbf{v}_{\text{femme}} \approx \mathbf{v}_{\text{reine}}
\]
\end{exemple}

Formellement, pour r\'esoudre l'analogie $a : b :: c : ?$, on cherche :
\[
d^* = \argmax_{d \in \mathcal{V} \setminus \{a,b,c\}} \cos(\mathbf{v}_d, \mathbf{v}_b - \mathbf{v}_a + \mathbf{v}_c)
\]

\begin{theoreme}[Structure lin\'eaire des analogies]
Sous l'hypoth\`ese que les plongements factorisent une matrice de PMI
(th\'eor\`eme~\ref{thm:neg-pmi}), les relations s\'emantiques r\'eguli\`eres
se manifestent comme des directions lin\'eaires dans l'espace de plongement :
\[
\mathbf{v}_b - \mathbf{v}_a \approx \mathbf{v}_{b'} - \mathbf{v}_{a'}
\]
pour toutes les paires $(a, b)$ et $(a', b')$ partageant la m\^eme relation.
\end{theoreme}

\subsection{Visualisation par t-SNE}

\begin{codeblock}[title=\textbf{Visualisation des plongements}]
\begin{minted}{python}
import numpy as np
from sklearn.manifold import TSNE
import matplotlib.pyplot as plt
from gensim.models import KeyedVectors

# Charger des plongements pre-entraines
wv = KeyedVectors.load_word2vec_format('GoogleNews-vectors.bin', binary=True)
words = ['king', 'queen', 'man', 'woman', 'prince', 'princess',
         'paris', 'france', 'berlin', 'germany', 'rome', 'italy']
vectors = np.array([wv[w] for w in words])

tsne = TSNE(n_components=2, random_state=42, perplexity=5)
coords = tsne.fit_transform(vectors)

plt.figure(figsize=(10, 8))
for i, word in enumerate(words):
    plt.scatter(coords[i, 0], coords[i, 1])
    plt.annotate(word, (coords[i, 0]+0.5, coords[i, 1]+0.5))
plt.title("Visualisation t-SNE des plongements Word2Vec")
plt.show()
\end{minted}
\end{codeblock}

% ------------------------------------------------------------
\section{\'Evaluation des plongements}
\label{sec:plg:evaluation}
% ------------------------------------------------------------

\subsection{\'Evaluation intrins\`eque}

\begin{itemize}
    \item \textbf{Analogies} : pr\'ecision sur des benchmarks d'analogies
          (Google analogy dataset, BATS).
    \item \textbf{Similarit\'e} : corr\'elation de Spearman entre la similarit\'e
          cosinus et les jugements humains (SimLex-999, WordSim-353).
\end{itemize}

\subsection{\'Evaluation extrins\`eque}

Performance des plongements comme features dans des t\^aches en aval :
classification de sentiments, NER, d\'esambiguation lexicale.

\begin{formules}[title=\textbf{M\'etriques d'\'evaluation des plongements}]
\begin{align}
\text{Pr\'ecision analogies} &= \frac{\text{analogies correctes}}{\text{total analogies}} \\
\rho_{\text{Spearman}} &= 1 - \frac{6 \sum_i d_i^2}{n(n^2 - 1)} & \text{(corr\'elation de rang)}
\end{align}
\end{formules}

% ------------------------------------------------------------
\section{Biais dans les plongements}
\label{sec:plg:biais}
% ------------------------------------------------------------

Les plongements appris sur de grands corpus textuels h\'eritent des biais
soci\'etaux pr\'esents dans les donn\'ees.

\begin{exemple}
Bolukbasi et al.\ (2016) ont montr\'e que :
\[
\mathbf{v}_{\text{informaticien}} - \mathbf{v}_{\text{homme}} + \mathbf{v}_{\text{femme}} \approx \mathbf{v}_{\text{m\'enag\`ere}}
\]
Ce r\'esultat refl\`ete des st\'er\'eotypes de genre pr\'esents dans les corpus.
\end{exemple}

\begin{definition}[D\'ebiaisage par projection]
Le d\'ebiaisage de Bolukbasi et al.\ consiste \`a :
\begin{enumerate}
    \item Identifier la direction de biais $\mathbf{g}$ (par ACP sur des paires
          de genre).
    \item Projeter les mots neutres orthogonalement \`a $\mathbf{g}$ :
          $\mathbf{v}_w' = \mathbf{v}_w - \inner{\mathbf{v}_w}{\mathbf{g}} \mathbf{g}$.
    \item \'Egaliser les paires d\'efinitionnelles.
\end{enumerate}
\end{definition}

% ------------------------------------------------------------
\section{Plongements contextuels vs.\ statiques}
\label{sec:plg:contextuels}
% ------------------------------------------------------------

Les plongements statiques (Word2Vec, GloVe, FastText) assignent un vecteur
unique \`a chaque mot, ind\'ependamment du contexte. Cela pose probl\`eme
pour les mots polysmiques :

\begin{exemple}
Le mot \og avocat \fg d\'esigne soit un fruit, soit un professionnel du
droit. Un plongement statique repr\'esente ces deux sens par un seul vecteur,
qui sera un compromis entre les deux.
\end{exemple}

Les plongements contextuels (ELMo, BERT, GPT) r\'esolvent ce probl\`eme en
produisant des repr\'esentations diff\'erentes selon le contexte. Ils seront
\'etudi\'es au chapitre~\ref{ch:pretraines}.

% ------------------------------------------------------------
\section{Impl\'ementation compl\`ete}
\label{sec:plg:implementation}
% ------------------------------------------------------------

\begin{codeblock}[title=\textbf{Entra\^inement Skip-gram avec PyTorch}]
\begin{minted}{python}
import torch
import torch.nn as nn

class SkipGram(nn.Module):
    def __init__(self, vocab_size, embed_dim):
        super().__init__()
        self.target_embed = nn.Embedding(vocab_size, embed_dim)
        self.context_embed = nn.Embedding(vocab_size, embed_dim)

    def forward(self, target, context, negatives):
        # target: (B,), context: (B,), negatives: (B, K)
        t = self.target_embed(target)          # (B, d)
        c = self.context_embed(context)        # (B, d)
        n = self.context_embed(negatives)      # (B, K, d)

        # Positive score
        pos_score = torch.sum(t * c, dim=1)    # (B,)
        pos_loss = -torch.log(torch.sigmoid(pos_score) + 1e-10)

        # Negative scores
        neg_score = torch.bmm(n, t.unsqueeze(2)).squeeze(2)  # (B, K)
        neg_loss = -torch.log(torch.sigmoid(-neg_score) + 1e-10).sum(dim=1)

        return (pos_loss + neg_loss).mean()

model = SkipGram(vocab_size=50000, embed_dim=300)
optimizer = torch.optim.Adam(model.parameters(), lr=0.001)
\end{minted}
\end{codeblock}

% ------------------------------------------------------------
\section{Exercices}
\label{sec:plg:exercices}
% ------------------------------------------------------------

\begin{exercice}[$\star$]
Montrez que la similarit\'e cosinus entre deux vecteurs one-hot est nulle
pour des mots diff\'erents.
\end{exercice}

\begin{exercice}[$\star\star$]
D\'erivez le gradient de l'objectif Skip-gram avec \'echantillonnage n\'egatif
par rapport aux vecteurs $\mathbf{u}_{w_i}$ et $\mathbf{v}_{w_o}$.
\end{exercice}

\begin{exercice}[$\star\star$]
Montrez que l'objectif GloVe est \'equivalent \`a une factorisation de matrice
pond\'er\'ee de la matrice $\log X_{ij}$.
\end{exercice}

\begin{exercice}[$\star\star$]
Entra\^inez des plongements Word2Vec (avec \texttt{gensim}) sur un corpus
fran\c{c}ais (Wikip\'edia). \'Evaluez les analogies ``roi -- homme + femme = ?''
et ``Paris -- France + Allemagne = ?''.
\end{exercice}

\begin{exercice}[$\star\star\star$]
Impl\'ementez la m\'ethode de d\'ebiaisage de Bolukbasi et al.\ (2016).
Mesurez le biais de genre avant et apr\`es correction sur un ensemble de
mots li\'es aux professions.
\end{exercice}
